% =====================================================================
%  P4.  A count is not an error bar: exact fluctuation floors for cell
%       means of arithmetic fields, with an application to
%       sum_{n<N} Lambda(n) mu(N-n).
%
%  Deployment version.  Written to deploy/SPEC.md.
%  Compile:  pdflatex P4-coherent-cell-floor.tex   (twice)
% =====================================================================
\documentclass[11pt,a4paper]{article}

\usepackage{amsmath}
\usepackage{amssymb}
\usepackage{amsthm}
\usepackage[margin=1in]{geometry}
\usepackage[colorlinks=true,linkcolor=blue,citecolor=blue]{hyperref}

\theoremstyle{plain}
\newtheorem{theorem}{Theorem}
\newtheorem{corollary}[theorem]{Corollary}
\newtheorem{proposition}[theorem]{Proposition}
\newtheorem{lemma}[theorem]{Lemma}
\theoremstyle{definition}
\newtheorem{observation}[theorem]{Observation}
\newtheorem{measurement}[theorem]{Measurement}
\theoremstyle{remark}
\newtheorem{note}[theorem]{Note}

\renewcommand{\SS}{\mathfrak{S}}
\newcommand{\AAA}{\mathfrak{A}}

% typographic slack: let TeX stretch a line rather than run into the margin
\emergencystretch=3em

\title{A count is not an error bar:\\
exact fluctuation floors for cell means of arithmetic fields}
\author{}
\date{}

\begin{document}
\maketitle

\begin{abstract}
Let $F(N)$ be an arithmetic field built by convolving a sign-valued
multiplicative function against a fixed nonnegative weight --- the
model case is $C(N)=\sum_{n<N}\Lambda(n)\mu(N-n)$, normalised as
$Z(N)=C(N)/\sqrt{V(N)}$. Partition a band of $N$ into cells indexed by
divisibility, and ask whether a cell mean differs from the band mean.
The usual error bar for that question is $\mathrm{sd}/\sqrt{n_c}$,
built from the cell count $n_c$.

That error bar is wrong, and not by a constant. We give the exact
fluctuation of any cell mean in closed form under the natural null in
which the signs are replaced by independent ones
(Lemma~\ref{lem:cellmom}): three convolution terms, no simulation. We
then show that this quantity does not decay like $n_c^{-1/2}$ but like
$(\log N)^{-1/2}$ (Observation~\ref{obs:coh}), because the summand
$u_c(v)$ is a \emph{coherent} sum rather than a self-averaging one:
every $N$ in the cell reads the same arithmetic input at the same
$v$. Over a factor $140$ in $N$, a count-based interval predicts the
error bar shrinks by $11.8$; it shrinks by $1.21$, against the $1.19$
that $(\log N)^{-1/2}$ predicts. An interval built from a count is
therefore about ten times too narrow at the top of that range relative
to the bottom, and in absolute terms the exact floor exceeds
$\mathrm{sd}(Z)/\sqrt{n_c}$ by factors of $5.8$ to $160$.

The exact floor is not merely a corrected noise level. Under the
label permutation that preserves every cell size while destroying the
correspondence between cells and arithmetic (Lemma~\ref{lem:placebo}),
the floor \emph{collapses} by factors from $3.8$ to $105$: it is
itself a second measurement of that correspondence, and quoting
against it is conservative by about two orders of magnitude. What
predicts its size is the excess of the shift's singular series over
same-cell pairs, a quantity that is scale-invariant by construction
(Proposition~\ref{prop:scaleinv}) --- so whatever produces the decay
of the effect it measures, it is not the arithmetic that produces its
size.

The conclusion applies to any cell mean of a field whose summands
share a common arithmetic input, and is independent of the Goldbach
context in which it was found.
\end{abstract}

% =====================================================================
\section{Introduction}

\subsection{The setting}

Let $N$ range over the even integers of a band, let
\begin{equation}\label{eq:field}
  C(N)=\sum_{n<N}\Lambda(n)\mu(N-n),
  \qquad
  V(N)=\sum_{v<N}\mu^2(v)\Lambda(N-v)^2,
  \qquad
  Z(N)=\frac{C(N)}{\sqrt{V(N)}},
\end{equation}
so that $Z$ has second moment $1$ under the null in which the $\mu(v)$
on their own support are replaced by independent signs; the exact
asymptotic $V(N)\sim\prod_{q\nmid N}(1-1/(q(q-1)))\,N\log N$ is
\cite[Prop.~1]{P3}, and the field \eqref{eq:field} is the scalar to
which the Huang--Li reduction of binary Goldbach \cite{HL} reduces.
Partition the
band into \emph{cells} indexed by \emph{depth} $d$, the number of
$3,5,7,11,13$ dividing $N$. The question is whether
$m_c-\overline m$, the difference between the cell mean of $Z$ and the
band mean, is larger than chance.

Answering it requires an error bar for $m_c-\overline m$. The reflex
is $\mathrm{sd}(Z)/\sqrt{n_c}$ with $n_c$ the cell count. This paper
shows that reflex fails here, gives the correct quantity in closed
form, and identifies the mechanism --- which is generic, not specific
to \eqref{eq:field}.

\subsection{Results}

\begin{itemize}
\item Lemma~\ref{lem:cellmom} gives $\mathrm{Var}(m_c-\overline m)$
exactly under the sign null, as three convolution terms.
\item Observation~\ref{obs:coh} identifies its decay: $(\log N)^{-1/2}$,
not $n_c^{-1/2}$.
\item Lemma~\ref{lem:placebo} is the control that separates a
property of the cell \emph{sizes} from a property of the
correspondence between cells and arithmetic.
\item Measurement~\ref{meas:placebo} runs it, and finds that the
effect survives while the floor itself does not --- so the floor is a
second measurement of the same correspondence.
\item Proposition~\ref{prop:scaleinv} shows the quantity predicting
the floor's size is scale-invariant.
\end{itemize}

\subsection{What is not claimed}\label{sec:notclaimed}

\begin{itemize}
\item \textbf{No decay exponent for the effect is quoted.} At the
three shallowest depths the amplitude does not clear the exact floor
at any scale measured, so nothing can be fitted there, and whether the
exponents differ by depth is not decided by the accessible range.
\item \textbf{No distributional law for $Z$ is asserted.}
\item \textbf{No consequence for the Goldbach problem.} The effect
measured here is lower order under every parameterisation considered
and does not threaten $C(N)=o(N)$.
\item \textbf{Observation~\ref{obs:coh} is offered for its form, not
its constant.} Its heuristic is quantitatively loose by a factor of
about $2.5$; the exponent is confirmed directly instead.
\end{itemize}

% =====================================================================
\section{The floor, in closed form}\label{sec:floor}

\begin{lemma}[exact cell moments]\label{lem:cellmom}
Let a band of even $N$ be partitioned into cells, let $c$ be a cell of
size $n_c$ inside a band of size $n$, let $m_c$ and $\overline m$ be
the means of $Z$ over $c$ and over the band, and set
\[
  u_c(v)\;=\;\sum_{N\in c}\frac{\Lambda(N-v)}{\sqrt{V(N)}},
  \qquad
  Q_{cd}\;=\;\sum_v\mu^2(v)\,u_c(v)\,u_d(v).
\]
Then, when the $\mu(v)$ on the surviving support are replaced by
independent signs,
\begin{equation}\label{eq:cellmom}
  \mathrm{Var}\bigl(m_c-\overline m\bigr)
  \;=\;\frac{Q_{cc}}{n_c^{2}}
  \;-\;\frac{2\,Q_{ca}}{n_c\,n}
  \;+\;\frac{Q_{aa}}{n^{2}},
\end{equation}
where $a$ denotes the whole band. Every term is computable exactly by
one convolution; no simulation is involved.
\end{lemma}

\begin{proof}
Under independent signs
$E[\varepsilon(v)\varepsilon(v')]=\mu^2(v)\delta_{vv'}$, so
$\mathrm{Cov}(n_cm_c,n_dm_d)=Q_{cd}$ exactly. Expanding
$\mathrm{Var}(m_c-\overline m)$ bilinearly gives the three terms.
\end{proof}

\begin{note}[all three terms are needed]\label{note:threeterms}
The substitution $u_c(v)\approx n_c/\sqrt V$, which gives the right
scale for $Q_{cc}/n_c^2$, is not specific to $c$; applied to all three
terms it returns the same value $S$ for each and hence $S-2S+S=0$. So
it cannot be used to estimate the variance of the difference. Exactly,
the variance is smaller than $Q_{cc}/n_c^2$ by a factor that is a
structural constant rather than noise.
\end{note}

\begin{measurement}[the closed form, checked against simulation]%
\label{meas:mc}
The lemma does not need simulation, but simulation is the only thing
that would catch an error in the formula, and every normalised cell
mean below is divided by it. Run in the band $(10^5,2\cdot10^5]$ with
$2000$ draws:
\[
\begin{array}{r|cccccc}
 \text{depth} & 0 & 1 & 2 & 3 & 4 & 5\\
 n_c & 19181 & 21097 & 8183 & 1421 & 115 & 3\\\hline
 \text{closed form} & 0.140080 & 0.052724 & 0.167579 & 0.273630
   & 0.378073 & 0.639330\\
 \text{MC}/\text{closed} & 0.9920 & 0.9993
   & 0.9918 & 1.0044 & 1.0157 & 1.0320
\end{array}
\]
--- worst deviation $0.0320$, at the depth whose cell holds three
elements, and the six ratios straddle $1$ three and three.

The six ratios are estimated from the \emph{same} draws, so their
deviations are strongly correlated and their common sign would be one
observation, not six; that is why the check is run at $2000$ draws
rather than at a number where the Monte-Carlo precision, here
$1/\sqrt{2\cdot1999}$, is comparable to the deviation being judged.

The structural constant of Note~\ref{note:threeterms} is visible
directly: at the top octave of $1.6\cdot10^7$,
$\mathrm{Var}/(Q_{cc}/n_c^2)$ is
$0.113119,\,0.016302,\,0.288567,\,0.551516$ at depths $0,1,3,5$, and
at a quarter of that $N$ the three shallow ratios agree to six digits
--- $0.113118,\,0.016303,\,0.288571$. Depth $5$ does not: it reads
$0.557654$ there against $0.551516$ at the top, differing in the third
digit, and it is much the smallest cell --- $266$ values of $N$ in the
top octave against $1\,687\,911$ at depth $1$
(\texttt{lab\_cellmom\_montecarlo.py}).
\end{measurement}

% =====================================================================
\section{The floor's uncertainty does not fall like a count}%
\label{sec:coh}

\begin{observation}[coherent sums]\label{obs:coh}
The error bar of Lemma~\ref{lem:cellmom} does \emph{not} decay like
$n_c^{-1/2}$. For fixed $v$, about $n_c/\log N$ of the terms of
$u_c(v)$ are nonzero, each of size $\log N/\sqrt V$, so
$u_c(v)\approx n_c/\sqrt V$ and
\[
  \frac{Q_{cc}}{n_c^2}
  \;\approx\;\frac{\sum_v\mu^2(v)}{V}
  \;\asymp\;\frac{1}{\log N},
\]
independently of $n_c$; and by Note~\ref{note:threeterms} the variance
itself is a fixed fraction of this. So the standard error falls like
$(\log N)^{-1/2}$.
\end{observation}

The heuristic is quantitatively loose and is offered for the
\emph{form} and not the constant: at $N=1.6\cdot10^7$,
\[
\begin{array}{r|cccccc}
 \text{depth} & 0 & 1 & 2 & 3 & 4 & 5\\\hline
 Q_{cc}/n_c^2 & 0.123577 & 0.121208 & 0.146026 & 0.184159 & 0.235171
   & 0.307074
\end{array}
\]
against the heuristic's
$\bigl(\sum_{v<N}\mu^2(v)\bigr)/V(N)=0.049540$. The row is not
monotone in the depth; its minimum is $0.121208$ at depth $1$.

\begin{measurement}[the exponent, fitted]\label{meas:exponent}
The form is confirmed directly. Fitting $\mathrm{se}\propto N^{-b}$ to
the exact floor across eight octaves gives
$b=0.039451,\,0.039671,\,0.039388$ at depths $2,1,0$, against the
$0.5$ that a count would give and against $0.036038$, the value of
$1/(2\langle\log N\rangle)$ with $\langle\log N\rangle$ the mean of
$\log N$ over the eight octave midpoints. Nothing turns on the third
decimal: the point is that $0.039$ is near $0.036$ and nowhere near
$0.5$ (\texttt{lab\_cell\_floor.py}).
\end{measurement}

\begin{note}[a count is not an error bar]\label{note:notanerrorbar}
The consequence is not confined to this field. Over a factor $140$ in
$N$ --- the sub-range of $10^5<N\le1.6\cdot10^7$ on which the exact
floor was fitted --- $n_c^{-1/2}$ says the error bar shrinks by
$11.8$; it shrinks by $1.21$, against the $1.19$ that
$(\log N)^{-1/2}$ predicts over the same factor. \textbf{An interval
built from a count is therefore about ten times too narrow at the top
of this range relative to the bottom}, and in absolute terms the
discrepancy is larger still: at the top octave the exact floor exceeds
$\mathrm{sd}(Z)/\sqrt{n_c}$ by factors of $5.8$ to $160$, growing with
the cell, exactly as $Q_{cc}/n_c^2\asymp1/\log N$ predicts. The cause
is that $u_c(v)$ is a coherent sum, not a self-averaging one, and the
same is true of any cell mean of a field whose summands share a common
arithmetic input.
\end{note}

% =====================================================================
\section{The control, and what running it shows}\label{sec:placebo}

\begin{lemma}[the placebo key]\label{lem:placebo}
Let the cells be indexed by a labelling $\ell(N)$, and let $\pi$ be a
permutation of the even integers in the band. Replacing $\ell$ by
$\ell\circ\pi$ preserves every cell size and leaves the field $Z$
pointwise unchanged, while destroying the correspondence between cells
and arithmetic. Any cell-indexed statistic that survives this
replacement is a property of the cell sizes and not of that
correspondence.
\end{lemma}

\begin{proof}
The permutation acts on labels only; the multiset
$\{Z(N):N\ \text{in the band}\}$ and the multiset of cell sizes are
both invariant, and the joint distribution of $(\ell\circ\pi,Z)$ is
that of independent labels.
\end{proof}

Against the exactly computed floor of Lemma~\ref{lem:cellmom}, the
cell means of the field \eqref{eq:field} are not zero. Over every
octave from $6.25\cdot10^4$ to $1.6\cdot10^7$, $\max_c|z_c|$ runs
$9.1$ to $13.0$ and clears Bonferroni in each; the signal is carried
by the deep cells, with depths $3,4,5$ at $z=-1.6,\,-4.5,\,-9.1$ at
the top octave while depths $0,1,2$ sit at $+0.0,\,+0.7,\,-0.1$.

\begin{measurement}[the effect survives the placebo; the floor does
not]\label{meas:placebo}
Run on the octave $(2\cdot10^6,\,4\cdot10^6]$, with cell sizes
preserved exactly and the floor recomputed from scratch for each
permutation:
\[
\begin{array}{r|cccccc}
 \text{depth} & 0 & 1 & 2 & 3 & 4 & 5\\
 n_c & 383617 & 421978 & 163568 & 28507 & 2263 & 67\\\hline
 z_c\ \text{(true)} & +0.1069 & +1.1133 & -0.2314 & -2.3990
   & -5.9997 & -11.0258
\end{array}
\]
against $\max_c|z_c|$ over ten label permutations reading
\[
  0.9359,\ 1.6073,\ 1.9921,\ 1.5687,\ 2.1392,\
  1.1216,\ 1.7178,\ 1.1348,\ 3.2040,\ 1.3192,
\]
mean $1.6741$, never above $3.3$. The
correlation between depth and $z_c$ is $-0.9106$ for the true
labelling and $+0.0042$ on average under permutation. So what is
detected is the correspondence between cells and divisibility, not the
cell sizes.

The same run corrects a natural reading of Lemma~\ref{lem:cellmom}.
One expects the floor $\mathrm{se}_c$ to be a property of the cell
sizes, hence nearly invariant under permutation. It is not: the floor
\textbf{collapses}, by factors from $3.8$ at depth $5$ to $105$ at
depth $0$ --- $1.2400\cdot10^{-1}$ against $1.1794\cdot10^{-3}$ there
(\texttt{lab\_mask\_placebo.py}).
\end{measurement}

\begin{note}[the floor is itself signal]\label{note:floorsignal}
The mechanism of that collapse is the three-term structure of
\eqref{eq:cellmom}. For a random subset of the band,
$u_c(v)\approx(n_c/n)\,u_a(v)$, and then
$Q_{cc}/n_c^2-2Q_{ca}/(n_cn)+Q_{aa}/n^2$ cancels to leading order,
leaving only the fluctuation around proportionality. An arithmetic
cell breaks that proportionality, and the floor is what is left over.
\textbf{The exact floor is therefore not a noise level but a second
measurement of the same correspondence}, and quoting $z_c$ against it
is conservative by about two orders of magnitude: with the permuted
floor in the denominator the depth-$5$ cell would read $z\approx-42$
rather than $-11$. The conservative figure is the one kept throughout.
\end{note}

% =====================================================================
\section{What predicts the floor's size}\label{sec:size}

Write $\SS_2(h)$ for the Hardy--Littlewood singular series of the
shift $h$ --- the local density of pairs $(p,p+h)$ --- and, for a cell
$c$ in the band, let $E_{\mathrm{same},c}[\SS_2]$ be the mean of
$\SS_2(N-N')$ over pairs $N,N'$ both lying in $c$, and
$E_{\mathrm{all}}[\SS_2]$ the same mean over all pairs in the band.
Put
\begin{equation}\label{eq:Dc}
  D_c\;:=\;E_{\mathrm{same},c}[\SS_2]-E_{\mathrm{all}}[\SS_2].
\end{equation}

\begin{measurement}[the floor tracks $D_c$, including its shape]%
\label{meas:Dc}
At the octave $(2\cdot10^6,\,4\cdot10^6]$:
\[
\begin{array}{r|cccccc}
 \text{depth} & 0 & 1 & 2 & 3 & 4 & 5\\\hline
 D_c & 0.302138 & 0.046766 & 0.412504 & 1.052129 & 1.946109
   & 3.171298\\
 \mathrm{se}_c & 0.12400 & 0.04662 & 0.14834 & 0.24178 & 0.33253
   & 0.43685
\end{array}
\]
with correlation $0.9805$ across depths. Neither row is monotone: both
dip at depth $1$, and that is the mechanism made visible. Depth $0$ is
the cell on which \emph{none} of $3,5,7,11,13$ divides $N$, and
excluding a class concentrates the shift just as fixing one does ---
if $3\nmid N$ and $3\nmid N'$ then both lie in $\{1,2\}\bmod3$ and
$3\mid h$ with probability $\tfrac12$ against $\tfrac13$ for a random
pair. Both ends of the depth index concentrate $h$; depth $1$ mixes
the patterns and washes out. Depth is therefore a coarse index for
$D_c$, and the floor follows the same coarse index
(\texttt{lab\_mask\_placebo.py}).
\end{measurement}

\begin{proposition}[the size mechanism is scale-invariant]%
\label{prop:scaleinv}
$D_c$ depends only on the distribution of the shift $h$ in the residue
classes mod $3,5,7,11,13$ that the cell fixes. That distribution is
the same at every scale, so $D_c$ is scale-invariant and predicts a
fitted exponent of zero at every depth.
\end{proposition}

\begin{proof}
For $N,N'$ in a band, $\SS_2(N-N')$ depends on $N-N'$ only through the
local factors at primes dividing $h=N-N'$, and the conditioning that
defines the cell fixes the residue of $N$ modulo each of
$3,5,7,11,13$. The induced law of $h$ modulo each of those primes is
therefore determined by the cell alone, and the band enters only
through the residues of $h$ at primes larger than $13$, whose joint
law over a band is the uniform one to $O(1/B)$, $B$ the band length.
Hence the
difference \eqref{eq:Dc} converges to a constant depending only on the
cell, and its scale exponent is $0$.
\end{proof}

\begin{measurement}[the exponent, measured]\label{meas:scaleinv}
Measured over $(10^6,2\cdot10^6]$, $(2\cdot10^6,4\cdot10^6]$ and
$(4\cdot10^6,8\cdot10^6]$, the fitted exponents are
\[
\begin{array}{r|cccccc}
 \text{depth} & 0 & 1 & 2 & 3 & 4 & 5\\\hline
 e & -0.000879 & -0.016226 & +0.000936 & +0.002904 & +0.000241
   & -0.008428
\end{array}
\]
--- zero to three decimals at every depth but one. The exception is
depth~$1$, and it is the depth at which $D_c=0.046766$ is six times
smaller than any other while the sampling error of the pair average is
the same $0.0013$ throughout; its fitted exponent flips sign from
$+0.065546$ to $-0.016226$ when the sample is taken ten times larger,
which is the signature of noise rather than of a trend. Depth~$1$ is
also the cell that carries no effect ($z=+1.11$ in
Measurement~\ref{meas:placebo}), so nothing rests on it
(\texttt{lab\_cell\_singular.py}).
\end{measurement}

Proposition~\ref{prop:scaleinv} settles one thing and deliberately not
another. Whatever produces the decay of the effect with $N$, it is not
the arithmetic excess that produces the floor's size, since that
excess does not decay at all. The decay exponent of the effect itself
is not determined here: at the three shallowest depths the amplitude
does not clear the exact floor at any scale measured, so no exponent
can be fitted there.

% =====================================================================
\section{Summary}

\begin{itemize}
\item Lemma~\ref{lem:cellmom}: the fluctuation of any cell mean under
the sign null is exactly
$Q_{cc}/n_c^2-2Q_{ca}/(n_cn)+Q_{aa}/n^2$; all three terms are needed
and all three are one convolution each.

\item Observation~\ref{obs:coh} and Note~\ref{note:notanerrorbar}:
that quantity decays like $(\log N)^{-1/2}$, not $n_c^{-1/2}$. A
count-based interval is about ten times too narrow across a factor
$140$ in $N$, and absolutely too narrow by factors of $5.8$ to $160$.

\item Lemma~\ref{lem:placebo} and Measurement~\ref{meas:placebo}: the
cell effect survives label permutation; the floor does not. The floor
is a second measurement of the same correspondence, and quoting
against it is conservative by two orders of magnitude.

\item Proposition~\ref{prop:scaleinv}: the quantity that predicts the
floor's size is scale-invariant, so it is not what makes the effect
decay.

\item The mechanism --- coherence of $u_c(v)$ across the cell --- is
generic. Any cell mean of a field whose summands share a common
arithmetic input has the same property.
\end{itemize}

% =====================================================================
\begin{thebibliography}{99}

\bibitem{HL} J.-J.~Huang and H.~Li, \emph{On the connection between
the Goldbach conjecture and the Elliott--Halberstam conjecture},
arXiv:2005.03811v2 [math.NT], 2022.

\bibitem{P3} \emph{The exact second moment of
$\sum_{n<N}\Lambda(n)\mu(N-n)$, and why Chowla's conjecture does not
control it}, companion paper.

\end{thebibliography}

% =====================================================================
\appendix
\section{Reproducibility}

\begin{center}
\begin{tabular}{lll}
\hline
Script & Result file & Supports\\
\hline
\texttt{lab\_cellmom\_montecarlo.py} &
 \texttt{lab\_cellmom\_montecarlo.txt} & Measurement~\ref{meas:mc}\\
\texttt{lab\_cell\_floor.py} & \texttt{lab\_cell\_floor.txt}
 & Measurement~\ref{meas:exponent}\\
\texttt{lab\_mask\_placebo.py} & \texttt{lab\_mask\_placebo.txt}
 & Measurements~\ref{meas:placebo},~\ref{meas:Dc}\\
\texttt{lab\_cell\_singular.py} & \texttt{lab\_cell\_singular.txt}
 & Measurement~\ref{meas:scaleinv}\\
\hline
\end{tabular}
\end{center}

Lemma~\ref{lem:cellmom}, Lemma~\ref{lem:placebo} and
Proposition~\ref{prop:scaleinv} are proved and use no computation.
The measurements state the band, the cell index and the statistic they
were taken of; each was run with its decision rule fixed in advance,
and each cell-indexed claim was run against the permutation control of
Lemma~\ref{lem:placebo}.

\end{document}
